% ---------------- 5.2 Baseline Algorithms (Non-A*) ----------------
% ---------------- 5.2 Ablation Studies ----------------
\subsection{Ablation Studies on A*}

To isolate the contribution of heuristic guidance, we evaluate our approach against three ablation baselines:

\begin{enumerate}
    \item \textbf{Cycle ablation:} We force classical cycle detection using Bellman--Ford, identifying negative cycles rather than optimizing for executable profitable paths.
    
    \item \textbf{Depth ablation:} We restrict search depth to shallow expansions. 
    The \texttt{simple\_1hop} variant evaluates only single outgoing actions, while 
    \texttt{simple\_2hop} evaluates minimal two-edge loops (the shortest possible cycle).
    
    \item \textbf{Heuristic ablation:} We remove heuristic guidance entirely by setting $h(n)=0$, reducing A* to Dijkstra’s algorithm and eliminating blockchain-aware signals such as congestion, liquidity, and order book depth.
\end{enumerate}


\subsubsection{Bellman--Ford Negative Cycle Detection}
\label{sec:bf-baseline}

We include Bellman--Ford negative-cycle detection as a classical arbitrage baseline. Given edge rates $r(e)$, we apply the log transform
\[
w(e) = -\log(r(e)),
\] % this cna be inline
so multiplicative rate products become additive path costs. Arbitrage opportunities correspond to negative-weight cycles. Bellman--Ford detects whether such a cycle exists in the transformed graph.

In our setting, Bellman--Ford is treated as a \emph{detection} method rather than an execution planner: it does not optimize for order-size-dependent slippage, transfer latency windows, or reliability constraints, and therefore frequently fails to produce an executable profitable cycle under real-world frictions.

\subsubsection{Greedy Spread-Based Strategy}
\label{sec:greedy-baseline}

We include shallow greedy baselines intended to emulate rapid spread scanning:

\begin{itemize}
    \item \textbf{1-hop edge scan (\texttt{simple\_1hop}):} tests whether any single outgoing action yields immediate positive gain. This is a minimal local check and typically cannot complete a full arbitrage cycle.
    \item \textbf{2-hop loop scan (\texttt{simple\_2hop}):} tests whether a two-action sequence can return to the original asset (a minimal loop-back).
\end{itemize}

These strategies are extremely fast, but they generally fail under realistic fees and transfer costs because profitable cross-exchange opportunities typically require coordinated trade+transfer actions.

\subsubsection{Dijkstra / Shortest Path Variant}
\label{sec:dijkstra-baseline}

We implement an unguided shortest-path baseline equivalent to A* with a zero heuristic:
\[
h(n) = 0.
\] % this thingy can be inline 
This reduces A* to Dijkstra’s algorithm, expanding nodes in non-decreasing accumulated cost $g(n)$. The baseline uses the same execution-aware edge model as our proposed methods (fees, withdrawals, gas, and transfer-time penalties encoded into $r(e)$), but provides no heuristic guidance.

Djiksrta represents the most accurate market costs as the paths are exact market data at time of search, but deviates under execution when accounting for market friction.

\begin{table}[t]
\centering
\small
\setlength{\tabcolsep}{4pt}
\caption{Baseline strategies used in evaluation.}
\label{tab:baseline-summary}
\begin{tabular}{ll}
\toprule
\textbf{Baseline} & \textbf{Description} \\
\midrule
Bellman--Ford & Negative-cycle detection on log-cost graph \\
Dijkstra ($h{=}0$) & Unguided cost-based search (A* with $h(n)=0$) \\
\texttt{2hop\_max} & Depth-limited unguided search (max 2 edges) \\
\texttt{simple\_1hop} & One-edge greedy scan (no cycle planning) \\
\texttt{simple\_2hop} & Two-edge loop-back greedy scan \\
\bottomrule
\end{tabular}
\end{table}


% ---------------- 5.5.1 Deterministic Heuristic Comparison ----------------
\subsubsection{Deterministic Heuristic Comparison}
\label{sec:deterministic-compare}

We compare heuristics and baselines on a synchronized live snapshot (timestamped run on 2026--02--12). The constructed graph contains 19 nodes, and we evaluate three start nodes: \texttt{binance:BUSD}, \texttt{binance:USDT}, and \texttt{kraken:USDT}, using 8 parallel workers. Table~\ref{tab:deterministic-compare} reports net profit and wall-clock time for each method at three order sizes.

\begin{table*}[h]
\centering
\small
\setlength{\tabcolsep}{6pt}
\caption{Heuristic-Guided A* Performance Across Order Sizes. 
Averages are computed over successful runs only. 
All profits are net after trading, withdrawal, and gas fees.}
\label{tab:heuristic-performance}
\begin{tabular}{lccccc}
\toprule
\textbf{Heuristic} 
& \textbf{Order Size} 
& \textbf{Avg Profit (\$)} 
& \textbf{Avg Profit (\%)} 
& \textbf{Avg Time (s)} 
& \textbf{Max Path Len} \\
\midrule

$h_1$
& \$100     & 0.04  & 0.040\% & 3.151 & 3 \\
& \$1,000   & 0.37  & 0.037\% & 0.002 & 3 \\
& \$10,000  & 3.69  & 0.037\% & 0.002 & 3 \\
\midrule

$h_2$
& \$100     & 0.04  & 0.040\% & 5.018 & 3 \\
& \$1,000   & 0.37  & 0.037\% & 0.002 & 3 \\
& \$10,000  & 3.69  & 0.037\% & 0.002 & 3 \\
\midrule

$h_3$
& \$100     & 0.04  & 0.040\% & 3.227 & 3 \\
& \$1,000   & 0.38  & 0.038\% & 0.007 & 3 \\
& \$10,000  & 3.76  & 0.038\% & 0.008 & 3 \\
\midrule

$h_4$
& \$100     & 0.04  & 0.040\% & 0.004 & 3 \\
& \$1,000   & 0.37  & 0.037\% & 0.004 & 3 \\
& \$10,000  & 3.69  & 0.037\% & 0.005 & 3 \\
\bottomrule
\end{tabular}
\end{table*}



Across this snapshot, \texttt{binance:BUSD} frequently yields no profitable path for both heuristics and baselines, while \texttt{binance:USDT} and \texttt{kraken:USDT} consistently produce profitable transfer+trade routes. Notably, Bellman--Ford fails to return a profitable cycle in these runs, reflecting the gap between classical negative-cycle detection and execution-aware feasibility under fees and operational constraints. Meanwhile, the shallow greedy scans (\texttt{simple\_1hop}, \texttt{simple\_2hop}) fail across all starts, indicating that profitable opportunities generally require multi-action planning rather than local edge inspection.
